\textbf{Phần 1. Trắc nghiệm lựa chọn}
\setcounterpageref{ex}{0}

\begin{ex} %[5X11V64]
Biểu đồ hình quạt tròn biểu diễn kết quả phân loại học lực của $120$ học sinh khối 5 gồm: Học sinh Xuất sắc chiếm $25\%$, học sinh Tốt chiếm $45\%$, học sinh Đạt chiếm $25\%$, còn lại là học sinh Cần cố gắng. Số học sinh Cần cố gắng của khối 5 là
\choice
{$4\text{ bạn}$}
{\True $6\text{ bạn}$}
{$8\text{ bạn}$}
{$10\text{ bạn}$}
\loigiai{
Tỉ số phần trăm học sinh Cần cố gắng là:
$$100\% - (25\% + 45\% + 25\%) = 5\%.$$
Số học sinh Cần cố gắng là:
$$120 \times 5\% = 6\text{ (bạn)}.$$
\textbf{Đáp án của câu hỏi là B}
}
\end{ex}

\begin{ex} %[5X11V64]
Một biểu đồ quạt tròn biểu diễn cơ cấu chi tiêu hàng tháng của một gia đình. Biết tiền ăn uống chiếm $40\%$, tiền học hành của con cái chiếm $25\%$, chi phí sinh hoạt khác chiếm $20\%$, còn lại là tiền tiết kiệm. Biết mỗi tháng gia đình tiết kiệm được $4\,500\,000\text{ đồng}$. Tổng thu nhập hàng tháng của gia đình đó là
\choice
{$25\,000\,000\text{ đồng}$}
{\True $30\,000\,000\text{ đồng}$}
{$32\,000\,000\text{ đồng}$}
{$28\,000\,000\text{ đồng}$}
\loigiai{
Tỉ số phần trăm tiền tiết kiệm là:
$$100\% - (40\% + 25\% + 20\%) = 15\%.$$
Tổng thu nhập hàng tháng của gia đình là:
$$4\,500\,000 : 15\% = 30\,000\,000\text{ (đồng)}.$$
\textbf{Đáp án của câu hỏi là B}
}
\end{ex}

\begin{ex} %[5X11V64]
Biểu đồ quạt tròn biểu diễn thể loại sách trong thư viện trường tiểu học: Sách giáo khoa chiếm $45\%$, Sách tham khảo chiếm $30\%$, Truyện tranh chiếm $15\%$, còn lại là Báo và tạp chí. Biết thư viện có tất cả $3\,600$ cuốn sách và báo các loại. Số lượng sách giáo khoa nhiều hơn số lượng truyện tranh là bao nhiêu cuốn?
\choice
{$900\text{ cuốn}$}
{\True $1\,080\text{ cuốn}$}
{$1\,200\text{ cuốn}$}
{$1\,260\text{ cuốn}$}
\loigiai{
Sách giáo khoa nhiều hơn truyện tranh số phần trăm là:
$$45\% - 15\% = 30\%.$$
Số lượng sách giáo khoa nhiều hơn truyện tranh là:
$$3\,600 \times 30\% = 1\,080\text{ (cuốn)}.$$
\textbf{Đáp án của câu hỏi là B}
}
\end{ex}

\begin{ex} %[5X11V64]
Trong một đợt khảo sát ý kiến của $500$ du khách về các điểm tham quan ở một thành phố, biểu đồ quạt tròn cho biết: $35\%$ thích tham quan di tích lịch sử, $30\%$ thích tham quan viện bảo tàng, $20\%$ thích công viên sinh thái, số còn lại thích phố đi bộ. Số du khách thích phố đi bộ là
\choice
{$60\text{ người}$}
{\True $75\text{ người}$}
{$80\text{ người}$}
{$90\text{ người}$}
\loigiai{
Tỉ số phần trăm du khách thích phố đi bộ là:
$$100\% - (35\% + 30\% + 20\%) = 15\%.$$
Số du khách thích phố đi bộ là:
$$500 \times 15\% = 75\text{ (người)}.$$
\textbf{Đáp án của câu hỏi là B}
}
\end{ex}

\begin{ex} %[5X11V64]
Biểu đồ hình quạt tròn biểu thị tỉ lệ học sinh tham gia các câu lạc bộ của trường: CLB Cờ vua chiếm $20\%$, CLB Bóng rổ chiếm $35\%$, CLB Âm nhạc chiếm $25\%$, CLB Mỹ thuật có $40$ học sinh. Tổng số học sinh tham gia cả $4$ câu lạc bộ là
\choice
{$180\text{ học sinh}$}
{\True $200\text{ học sinh}$}
{$240\text{ học sinh}$}
{$250\text{ học sinh}$}
\loigiai{
Tỉ số phần trăm học sinh tham gia CLB Mỹ thuật là:
$$100\% - (20\% + 35\% + 25\%) = 20\%.$$
Tổng số học sinh tham gia cả $4$ CLB là:
$$40 : 20\% = 200\text{ (học sinh)}.$$
\textbf{Đáp án của câu hỏi là B}
}
\end{ex}

\begin{ex} %[5X11V64]
Trên biểu đồ quạt tròn biểu diễn diện tích trồng trọt một trang trại rộng $12\text{ ha}$: Diện tích trồng lúa chiếm $50\%$, diện tích cây ăn quả chiếm $30\%$, diện tích trồng hoa chiếm $15\%$, diện tích hồ thả cá chiếm $5\%$. Diện tích trồng hoa của trang trại đó là
\choice
{$1{,}5\text{ ha}$}
{\True $1{,}8\text{ ha}$}
{$2\text{ ha}$}
{$1{,}2\text{ ha}$}
\loigiai{
Diện tích trồng hoa là:
$$12 \times 15\% = 1{,}8\text{ (ha)}.$$
\textbf{Đáp án của câu hỏi là B}
}
\end{ex}

\begin{ex} %[5X11V64]
Một biểu đồ quạt tròn biểu diễn cơ cấu các loại xe gửi trong một bãi đỗ xe: Xe máy chiếm $65\%$, Xe đạp chiếm $15\%$, Ô tô chiếm $20\%$. Biết có $260$ chiếc xe máy trong bãi. Số lượng xe ô tô gửi trong bãi là
\choice
{$60\text{ chiếc}$}
{\True $80\text{ chiếc}$}
{$70\text{ chiếc}$}
{$90\text{ chiếc}$}
\loigiai{
Tổng số xe gửi trong bãi là:
$$260 : 65\% = 400\text{ (chiếc)}.$$
Số lượng xe ô tô gửi trong bãi là:
$$400 \times 20\% = 80\text{ (chiếc)}.$$
\textbf{Đáp án của câu hỏi là B}
}
\end{ex}

\begin{ex} %[5X11V64]
Biểu đồ quạt tròn cho biết cơ cấu khách hàng mua hàng tại một siêu thị trong ngày: Khách thanh toán tiền mặt chiếm $30\%$, thanh toán qua ví điện tử chiếm $45\%$, thanh toán qua thẻ ngân hàng chiếm $25\%$. Nếu có $360$ khách thanh toán bằng ví điện tử thì có bao nhiêu khách thanh toán bằng tiền mặt?
\choice
{$200\text{ khách}$}
{\True $240\text{ khách}$}
{$220\text{ khách}$}
{$250\text{ khách}$}
\loigiai{
Tổng số khách hàng đến siêu thị trong ngày là:
$$360 : 45\% = 800\text{ (khách)}.$$
Số khách thanh toán bằng tiền mặt là:
$$800 \times 30\% = 240\text{ (khách)}.$$
\textbf{Đáp án của câu hỏi là B}
}
\end{ex}

\begin{ex} %[5X11V64]
Quan sát biểu đồ quạt tròn thấy phần biểu diễn của một nhóm thành phần là hình quạt có góc ở tâm là $90^{\circ}$. Tỉ số phần trăm của nhóm thành phần đó so với toàn bộ hình tròn là
\choice
{$20\%$}
{\True $25\%$}
{$30\%$}
{$33{,}33\%$}
\loigiai{
Cả hình tròn có góc ở tâm bằng $360^{\circ}$ ứng với $100\%$.\\
Góc $90^{\circ}$ chiếm tỉ lệ:
$$\dfrac{90^{\circ}}{360^{\circ}} = \dfrac{1}{4} = 25\%.$$
\textbf{Đáp án của câu hỏi là B}
}
\end{ex}

\begin{ex} %[5X11V64]
Biểu đồ quạt tròn biểu diễn kết quả bình chọn nhân vật lịch sử yêu thích của $240$ học sinh: Nhân vật A chiếm $40\%$, Nhân vật B chiếm $35\%$, Nhân vật C chiếm $25\%$. Số phiếu bình chọn cho Nhân vật A nhiều hơn Nhân vật C là
\choice
{$30\text{ phiếu}$}
{\True $36\text{ phiếu}$}
{$40\text{ phiếu}$}
{$42\text{ phiếu}$}
\loigiai{
Hiệu tỉ số phần trăm giữa Nhân vật A và Nhân vật C là:
$$40\% - 25\% = 15\%.$$
Số phiếu bình chọn Nhân vật A nhiều hơn Nhân vật C là:
$$240 \times 15\% = 36\text{ (phiếu)}.$$
\textbf{Đáp án của câu hỏi là B}
}
\end{ex}

\begin{ex} %[5X11V64]
Một cửa hàng bán ba loại nước ép hoa quả gồm: Nước cam, nước táo và nước dưa hấu. Biểu đồ quạt tròn cho biết: Nước cam chiếm $45\%$, nước táo chiếm $35\%$. Biết trong ngày cửa hàng bán được $30$ cốc nước dưa hấu. Tổng số cốc nước ép cửa hàng đã bán trong ngày là
\choice
{$120\text{ cốc}$}
{\True $150\text{ cốc}$}
{$160\text{ cốc}$}
{$180\text{ cốc}$}
\loigiai{
Tỉ số phần trăm nước dưa hấu bán được là:
$$100\% - (45\% + 35\%) = 20\%.$$
Tổng số cốc nước ép đã bán là:
$$30 : 20\% = 150\text{ (cốc)}.$$
\textbf{Đáp án của câu hỏi là B}
}
\end{ex}

\begin{ex} %[5X11V64]
Biểu đồ quạt tròn biểu diễn tỉ lệ yêu thích các loại quả của $80$ học sinh: Cam ($25\%$), Xoài ($40\%$), Táo ($20\%$), Chuối ($15\%$). Số học sinh thích xoài gấp mấy lần số học sinh thích táo?
\choice
{$1{,}5\text{ lần}$}
{\True $2\text{ lần}$}
{$2{,}5\text{ lần}$}
{$3\text{ lần}$}
\loigiai{
Tỉ số giữa tỉ lệ thích xoài và thích táo là:
$$40\% : 20\% = 2\text{ (lần)}.$$
Do đó số bạn thích xoài gấp $2$ lần số bạn thích táo.\\
\textbf{Đáp án của câu hỏi là B}
}
\end{ex}

\textbf{Phần 2. Trắc nghiệm đúng sai}
\setcounterpageref{ex}{0}

\begin{ex} %[5X11V64]
Một biểu đồ hình quạt tròn biểu diễn cơ cấu các khoản chi trong chuyến dã ngoại của lớp 5A gồm: Tiền xe đi lại chiếm $35\%$, tiền ăn uống chiếm $45\%$, vé tham quan chiếm $15\%$, tiền nước uống dự phòng chiếm $5\%$. Biết tổng chi phí của chuyến đi là $6\,000\,000\text{ đồng}$. Xét tính đúng/sai của các nhận định sau:
\choiceTF
{\True Chi phí ăn uống là khoản chi tốn kém nhất, hết $2\,700\,000\text{ đồng}$}
{\True Chi phí tiền xe đi lại là $2\,100\,000\text{ đồng}$}
{Chi phí vé tham quan hết $1\,200\,000\text{ đồng}$}
{\True Chi phí tiền nước uống dự phòng hết $300\,000\text{ đồng}$}
\loigiai{
\begin{itemchoice}
\item Đúng. Ăn uống: $6\,000\,000 \times 45\% = 2\,700\,000\text{ đồng}$.
\item Đúng. Tiền xe: $6\,000\,000 \times 35\% = 2\,100\,000\text{ đồng}$.
\item Sai. Vé tham quan: $6\,000\,000 \times 15\% = 900\,000\text{ đồng}$ (chứ không phải $1\,200\,000\text{ đồng}$).
\item Đúng. Nước uống: $6\,000\,000 \times 5\% = 300\,000\text{ đồng}$.
\end{itemchoice}
}
\end{ex}

\begin{ex} %[5X11V64]
Biểu đồ quạt tròn biểu diễn kết quả khảo sát phương tiện đi làm của $200$ nhân viên trong một công ty: Đi xe máy chiếm $60\%$, đi xe buýt chiếm $25\%$, đi ô tô cá nhân chiếm $10\%$, đi xe đạp chiếm $5\%$. Xét tính đúng/sai của các nhận định sau:
\choiceTF
{\True Có $120$ nhân viên đi làm bằng xe máy}
{\True Số nhân viên đi xe buýt nhiều hơn số nhân viên đi ô tô cá nhân là $30$ người}
{Số nhân viên đi xe đạp là $15$ người}
{\True Tỉ số giữa số người đi xe buýt và số người đi xe máy là $\dfrac{5}{12}$}
\loigiai{
\begin{itemchoice}
\item Đúng. $200 \times 60\% = 120$ (người).
\item Đúng. Xe buýt ($200 \times 25\% = 50$), ô tô ($200 \times 10\% = 20$) $\Rightarrow 50 - 20 = 30$ (người).
\item Sai. Đi xe đạp là $200 \times 5\% = 10$ người (không phải 15 người).
\item Đúng. Tỉ số: $25\% : 60\% = \dfrac{25}{60} = \dfrac{5}{12}$.
\end{itemchoice}
}
\end{ex}

\textbf{Phần 3. Câu trả lời ngắn}
\setcounterpageref{ex}{0}

\begin{ex} %[5X11V64]
Biểu đồ quạt tròn biểu diễn các loại hoa trong vườn trường: Hoa hồng chiếm $40\%$, hoa cúc chiếm $35\%$, còn lại là $50$ cây hoa đồng tiền. Hỏi vườn trường có tất cả bao nhiêu cây hoa?
\shortans{200}
\loigiai{
Tỉ số phần trăm hoa đồng tiền là:
$$100\% - (40\% + 35\%) = 25\%.$$
Tổng số cây hoa trong vườn trường là:
$$50 : 25\% = 200\text{ (cây)}.$$
}
\end{ex}

\begin{ex} %[5X11V64]
Một biểu đồ quạt tròn cho biết tỉ lệ bình chọn tiết mục văn nghệ xuất sắc nhất: Tiết mục Hát chiếm $48\%$, Tiết mục Múa chiếm $32\%$, Tiết mục Kịch chiếm $20\%$. Biết có $75$ phiếu bình chọn cho Tiết mục Kịch. Hỏi Tiết mục Hát nhận được bao nhiêu phiếu bình chọn?
\shortans{180}
\loigiai{
Tổng số phiếu bình chọn là:
$$75 : 20\% = 375\text{ (phiếu)}.$$
Số phiếu bình chọn cho Tiết mục Hát là:
$$375 \times 48\% = 180\text{ (phiếu)}.$$
}
\end{ex}

\begin{ex} %[5X11V64]
Biểu đồ quạt tròn biểu diễn sản lượng thu hoạch ngũ cốc của một hợp tác xã gồm: Thóc chiếm $60\%$, Ngô chiếm $25\%$, Khoai chiếm $15\%$. Biết sản lượng thóc thu hoạch được là $180\text{ tấn}$. Hỏi sản lượng ngô thu hoạch được là bao nhiêu tấn?
\shortans{75}
\loigiai{
Tổng sản lượng ngũ cốc là:
$$180 : 60\% = 300\text{ (tấn)}.$$
Sản lượng ngô thu hoạch được là:
$$300 \times 25\% = 75\text{ (tấn)}.$$
}
\end{ex}

\begin{ex} %[5X11V64]
Biểu đồ quạt tròn thể hiện sự yêu thích các môn học của $150$ học sinh: Toán ($36\%$), Tiếng Việt ($28\%$), Tiếng Anh ($24\%$), Lịch sử và Địa lí ($12\%$). Hỏi số học sinh thích môn Toán nhiều hơn số học sinh thích môn Tiếng Anh là bao nhiêu bạn?
\shortans{18}
\loigiai{
Tỉ số phần trăm học sinh thích Toán hơn Tiếng Anh là:
$$36\% - 24\% = 12\%.$$
Số học sinh chênh lệch là:
$$150 \times 12\% = 18\text{ (bạn)}.$$
}
\end{ex}

\begin{ex} %[5X11V64]
Một nông trại nuôi $600$ con gia cầm. Biểu đồ quạt tròn cho biết: Gà chiếm $55\%$, Vịt chiếm $30\%$, Ngan chiếm $15\%$. Hỏi nông trại đó nuôi bao nhiêu con vịt?
\shortans{180}
\loigiai{
Số con vịt nông trại nuôi là:
$$600 \times 30\% = 180\text{ (con)}.$$
}
\end{ex}

\begin{ex} %[5X11V64]
Một người chia đều thời gian $24$ giờ trong ngày theo biểu đồ quạt: Học tập chiếm $25\%$, Ngủ chiếm $37{,}5\%$, Thể thao và giải trí chiếm $12{,}5\%$, các hoạt động cá nhân khác chiếm $25\%$. Hỏi trong một ngày người đó dành bao nhiêu giờ để ngủ?
\shortans{9}
\loigiai{
Số giờ người đó dành để ngủ trong một ngày là:
$$24 \times 37{,}5\% = 24 \times 0{,}375 = 9\text{ (giờ)}.$$
}
\end{ex}
